\documentclass[12pt,a4paper]{report}

% =========================================================
% BIG DATA SYSTEMS (24CI51) - PROJECT / SELF-LEARNING REPORT
% =========================================================

% ---------------- Packages ----------------
\usepackage[a4paper,top=1.55in,bottom=1in,left=1.25in,right=1in,
            headheight=75pt,headsep=18pt]{geometry}
\usepackage{setspace}
\usepackage{times}
\usepackage{graphicx}
\usepackage{float}
\usepackage{array}
\usepackage{tabularx}
\usepackage{longtable}
\usepackage{titlesec}
\usepackage{tocloft}
\usepackage{fancyhdr}
\usepackage{hyperref}
\usepackage{ragged2e}
\usepackage{enumitem}
\usepackage{caption}
\usepackage{booktabs}
\usepackage{multirow}
\usepackage{xcolor}

\hypersetup{
    colorlinks=true,
    linkcolor=black,
    urlcolor=black,
    citecolor=black
}

\onehalfspacing
\setlength{\parindent}{0.5in}
\setlength{\parskip}{0pt}
\justifying
\setlist[itemize]{leftmargin=*}
\setlist[enumerate]{leftmargin=*}
\renewcommand{\arraystretch}{1.2}

% ---------------- Project Details ----------------
\newcommand{\ProjectTitle}{BIG DATA SYSTEMS: TECHNOLOGY STUDY AND IMPLEMENTATION}
\newcommand{\CourseCode}{24CI51 -- BIG DATA SYSTEMS}
\newcommand{\AcademicYear}{2025--2026}
\newcommand{\Department}{DEPARTMENT OF CSE (ARTIFICIAL INTELLIGENCE \& MACHINE LEARNING)}
\newcommand{\StudentOne}{[NAME 1]}
\newcommand{\USNOne}{[USN 1]}
\newcommand{\GuideOne}{Mr. Subash N}
\newcommand{\GuideDesignation}{Assistant Professor}
\newcommand{\HOD}{Dr. Siddesh G M}
\newcommand{\HODDesignation}{Professor \& Head}

% ---------------- Header/Footer ----------------
% Keep the institute header image at images/rit_header.png.
% If the image is unavailable, comment out \RamaiahHeaderImage in the header definitions.
\newcommand{\RamaiahHeaderImage}{%
  \includegraphics[width=1\textwidth,height=4cm,keepaspectratio]{images/rit_header.png}%
}

\pagestyle{fancy}
\fancyhf{}
\fancyhead[C]{\RamaiahHeaderImage}
\fancyfoot[L]{\small \AcademicYear}
\fancyfoot[C]{\small Big Data Systems (24CI51)}
\fancyfoot[R]{\small \thepage}
\renewcommand{\headrulewidth}{0pt}
\renewcommand{\footrulewidth}{0.4pt}

\fancypagestyle{plain}{%
  \fancyhf{}
  \fancyhead[C]{\RamaiahHeaderImage}
  \fancyfoot[L]{\small \AcademicYear}
  \fancyfoot[C]{\small Big Data Systems (24CI51)}
  \fancyfoot[R]{\small \thepage}
  \renewcommand{\headrulewidth}{0pt}
  \renewcommand{\footrulewidth}{0.4pt}
}

\titleformat{\chapter}[display]
  {\normalfont\bfseries\centering\Large}
  {\chaptername\ \thechapter}{10pt}{\Large}
\titleformat{\section}{\normalfont\bfseries\large}{\thesection}{1em}{}
\titleformat{\subsection}{\normalfont\bfseries}{\thesubsection}{1em}{}

% ---------------- Reusable report helpers ----------------
\newcommand{\placeholder}[1]{\textit{[Enter #1 here.]}}
\newcommand{\activityheading}[2]{%
  \section{Activity #1: #2}%
}
\newcommand{\insertfigure}[3][0.85\textwidth]{%
\begin{figure}[H]
    \centering
    % Replace the filename below with the relevant image.
    \includegraphics[width=#1]{#2}
    \caption{#3}
    \label{fig:#2}
\end{figure}
}

\begin{document}

% =========================================================
% TITLE PAGE
% =========================================================
\begin{titlepage}
\centering
\thispagestyle{empty}

\vspace*{0.4cm}
{\bfseries\Large A REPORT ON\\[0.3cm]}
{\bfseries\LARGE \ProjectTitle\\[0.6cm]}

{\bfseries for the course\\[0.2cm]}
{\bfseries \CourseCode\\[0.8cm]}

Submitted in partial fulfillment of the academic requirements
for the course work of\\[0.2cm]
{\bfseries Bachelor of Engineering\\}
in\\
{\bfseries Computer Science and Engineering\\[0.2cm]}
{\bfseries (Artificial Intelligence \& Machine Learning)\\[0.8cm]}

{\bfseries Submitted by:\\[0.2cm]}
\StudentOne\ -- \USNOne\\
[0.8cm]

{\bfseries Under the Guidance of:\\[0.2cm]}
\GuideOne\\
\GuideDesignation\\[1.0cm]

{\bfseries M. S. RAMAIAH INSTITUTE OF TECHNOLOGY\\}
{\bfseries BENGALURU -- 560054\\}
{\bfseries \AcademicYear}

\end{titlepage}

\pagenumbering{roman}
\setcounter{page}{1}

% =========================================================
% CERTIFICATE
% =========================================================
\chapter*{CERTIFICATE}
\addcontentsline{toc}{chapter}{Certificate}

Certified that the report entitled ``\ProjectTitle'' submitted by
\StudentOne\ -- \USNOne, \StudentTwo\ -- \USNTwo, and
\StudentThree\ -- \USNThree, bonafide students of M. S. Ramaiah
Institute of Technology, Bengaluru, is a record of the work carried
out by them for the course \CourseCode\ during the academic year
\AcademicYear. The report is submitted in partial fulfillment of the
academic requirements prescribed by the institute.

\vspace{2.0cm}

\begin{tabularx}{\textwidth}{X X}
\centering Signature of the Guide &
\centering\arraybackslash Signature of the HoD\\[0.4cm]
\centering \GuideOne &
\centering\arraybackslash \HOD\\
\centering \GuideDesignation &
\centering\arraybackslash \HODDesignation\\
\end{tabularx}

\clearpage

% =========================================================
% DECLARATION
% =========================================================
\chapter*{DECLARATION}
\addcontentsline{toc}{chapter}{Declaration}

We hereby declare that the report entitled ``\ProjectTitle'' is the
record of work carried out by us at M. S. Ramaiah Institute of
Technology, Bengaluru, as part of the course \CourseCode. The work
presented in this report is prepared by us and has not been submitted
in part or full for the award of any other degree or diploma.

\vspace{2.0cm}

\begin{tabularx}{\textwidth}{X X}
Signature & Signature\\[0.4cm]
\StudentOne & \StudentTwo\\
\USNOne & \USNTwo\\[1.0cm]
\multicolumn{2}{l}{Signature:}\\[0.4cm]

\end{tabularx}

\clearpage

% =========================================================
% ACKNOWLEDGEMENT
% =========================================================
\chapter*{ACKNOWLEDGEMENT}
\addcontentsline{toc}{chapter}{Acknowledgement}

We express our sincere gratitude to the management and Principal of
M. S. Ramaiah Institute of Technology for providing the facilities
and support required to complete this work. We thank
\HOD, Professor \& Head, Department of CSE (Artificial Intelligence
\& Machine Learning), for the encouragement and guidance provided.

We express our sincere thanks to our course guide, \GuideOne,
\GuideDesignation, for valuable guidance and suggestions throughout
the completion of this work. We also thank the faculty members,
technical staff, friends, and family for their support.

\vspace{1cm}
\begin{flushright}
\StudentOne\ -- \USNOne\\

\end{flushright}

\clearpage

% =========================================================
% ABSTRACT
% =========================================================
\chapter*{ABSTRACT}
\addcontentsline{toc}{chapter}{Abstract}

\placeholder{an abstract describing the purpose of the Big Data Systems
activities, the five technologies studied, the implementation approach,
the datasets or simulations used, major observations, and conclusions.
Mention that synthetic or sample data is used wherever applicable.}

\clearpage

\tableofcontents
\clearpage
\listoffigures
\addcontentsline{toc}{chapter}{List of Figures}
\clearpage
\listoftables
\addcontentsline{toc}{chapter}{List of Tables}
\clearpage

\pagenumbering{arabic}

% =========================================================
% CHAPTER 1 - INTRODUCTION
% =========================================================
\chapter{INTRODUCTION}

\section{Introduction to Big Data Systems}
\placeholder{Introduce Big Data, its characteristics, the need for
distributed storage and processing, and the role of batch processing,
stream processing, NoSQL databases, and cloud analytics.}

\section{Problem Definition}
\placeholder{Describe the real-world data-management and analytics
problems addressed by the five activities: sales analysis, traffic
monitoring, patient vital-sign stream monitoring, product-catalog
management, and application-log analytics.}

\section{Objectives}
\begin{itemize}
    \item Study the architecture, concepts, and applications of Apache Hive.
    \item Use Apache Kafka concepts to demonstrate continuous data ingestion.
    \item Explore Apache Flink for real-time stream processing and alert generation.
    \item Implement MongoDB document storage, retrieval, updates, and queries.
    \item Explore a cloud-based platform for large-scale log analytics.
    \item Analyze outputs and document observations, limitations, and conclusions.
\end{itemize}

\section{Scope of the Work}
\placeholder{Describe the scope of the five activities. State the datasets,
tools, simulations, and environments used. Mention that simulated or
synthetic data is used where real operational data is unavailable.}

\section{Overview of Activities}
\begin{table}[H]
\centering
\caption{Overview of Big Data Systems Activities}
\begin{tabularx}{\textwidth}{|c|l|X|}
\hline
\textbf{No.} & \textbf{Technology} & \textbf{Application Scenario}\\
\hline
1 & Apache Hive & Supermarket sales analysis and inventory insights\\
\hline
2 & Apache Kafka & Smart-city traffic data ingestion and congestion alerts\\
\hline
3 & Apache Flink & Real-time processing of synthetic patient vital-sign data\\
\hline
4 & MongoDB & E-commerce product-catalog storage and retrieval\\
\hline
5 & Cloud-Based Big Data Analytics & Application and web-server log analysis\\
\hline
\end{tabularx}
\end{table}

% =========================================================
% CHAPTER 2 - REQUIREMENT SPECIFICATIONS
% =========================================================
\chapter{PROJECT REQUIREMENT SPECIFICATIONS}

\section{Software Requirements}
\begin{table}[H]
\centering
\caption{Software Requirements}
\begin{tabularx}{\textwidth}{|l|X|}
\hline
\textbf{Component} & \textbf{Technology / Requirement}\\
\hline
Operating System & [Linux / Windows and version]\\
\hline
Programming Language & [Python / Java / other language used]\\
\hline
Data Warehousing & Apache Hive and HiveQL\\
\hline
Message Streaming & Apache Kafka\\
\hline
Stream Processing & Apache Flink\\
\hline
NoSQL Database & MongoDB\\
\hline
Cloud Analytics & One selected platform: AWS, Microsoft Azure, or Google Cloud\\
\hline
Development Tools & [IDE, terminal, notebook, or other tools]\\
\hline
Visualization & [Tool/library used, if applicable]\\
\hline
\end{tabularx}
\end{table}

\section{Hardware Requirements}
\begin{table}[H]
\centering
\caption{Hardware Requirements}
\begin{tabularx}{\textwidth}{|l|X|}
\hline
\textbf{Component} & \textbf{Minimum / Available Specification}\\
\hline
Processor & [Processor model and speed]\\
\hline
RAM & [Available RAM]\\
\hline
Storage & [Available disk space]\\
\hline
Network & Internet or local network connectivity, as required\\
\hline
Additional Hardware & [Mention additional devices, if used]\\
\hline
\end{tabularx}
\end{table}

\section{Functional Requirements}
\begin{itemize}
    \item Create Hive tables and execute queries on sales transaction data.
    \item Publish traffic sensor records to Kafka topics and consume incoming events.
    \item Process a stream of synthetic patient readings using Flink and identify values outside demonstration thresholds.
    \item Insert, retrieve, update, search, and delete product documents in MongoDB.
    \item Upload or access sample server logs using a selected cloud platform and analyze log patterns.
    \item Generate outputs, alerts, summaries, tables, or visualizations appropriate to each activity.
\end{itemize}

\section{Non-Functional Requirements}
\begin{itemize}
    \item \textbf{Performance:} The solution should process the selected sample data within reasonable time.
    \item \textbf{Scalability:} The technologies studied should support approaches for handling growing data volumes.
    \item \textbf{Reliability:} The implementation should produce consistent outputs for the same input data.
    \item \textbf{Usability:} Queries, scripts, and results should be documented clearly.
    \item \textbf{Security:} Access credentials and sensitive information should not be exposed in the report.
    \item \textbf{Cost Awareness:} Cloud usage and estimated resource costs should be documented where applicable.
\end{itemize}

\section{Data Requirements}
\begin{table}[H]
\centering
\caption{Datasets Required for the Activities}
\begin{tabularx}{\textwidth}{|l|X|X|}
\hline
\textbf{Activity} & \textbf{Dataset / Data Fields} & \textbf{Purpose}\\
\hline
Hive & Sales records: transaction ID, date, product, quantity, price, branch & Revenue, product, branch, and monthly sales analysis\\
\hline
Kafka & Traffic events: timestamp, junction ID, vehicle count, speed & Continuous ingestion and congestion detection\\
\hline
Flink & Synthetic readings: timestamp, patient/device ID, heart rate, oxygen saturation, temperature & Demonstration of stream processing and threshold alerts\\
\hline
MongoDB & Product name, category, price, brand, specifications, rating, availability & Product catalog and CRUD operations\\
\hline
Cloud Analytics & Sample application/web-server logs, timestamps, status codes, request details & Error frequency, request volume, spikes, and usage patterns\\
\hline
\end{tabularx}
\end{table}

% =========================================================
% CHAPTER 3 - TECHNOLOGY STUDY AND DESIGN
% =========================================================
\chapter{TECHNOLOGY STUDY AND SYSTEM DESIGN}

\section{Overall System Design}
\placeholder{Describe how the five activities are organized and how input
data is collected, stored, processed, analyzed, and presented. The activities
may be implemented independently.}

\begin{figure}[H]
    \centering
    % Add the overall architecture diagram in the images folder.
    % \includegraphics[width=0.90\textwidth]{images/overall_architecture.png}
    \fbox{\parbox[c][5cm][c]{0.85\textwidth}{\centering
    Insert Overall Big Data Systems Architecture Diagram}}
    \caption{Overall Big Data Systems Architecture}
    \label{fig:overall_architecture}
\end{figure}

\section{Apache Hive}
\subsection{Introduction and Architecture}
\placeholder{Explain Apache Hive, its role in data warehousing, and its
architecture. Include relevant components such as the user interface,
driver, compiler, execution engine, metastore, and storage layer, as
applicable to the environment used.}

\subsection{HiveQL and Data Warehousing Concepts}
\placeholder{Explain HiveQL, databases, tables, schemas, data types,
partitions or buckets if used, and the difference between Hive and a
traditional relational database.}

\subsection{Proposed Hive-Based Sales Analysis}
\placeholder{Describe how supermarket transaction records are stored in
Hive and how the data is queried to identify total revenue, best-selling
products, peak sales hours, branch-wise revenue, and monthly trends.}

\section{Apache Kafka}
\subsection{Introduction and Architecture}
\placeholder{Explain Kafka and the roles of brokers, topics, partitions,
producers, consumers, and consumer groups, as applicable.}

\subsection{Producer-Consumer Model}
\placeholder{Describe how traffic sensor events are generated, published
to a Kafka topic, consumed by an application, and processed.}

\subsection{Proposed Traffic Monitoring Design}
\placeholder{Explain the traffic-event fields and the method used to
identify congestion using predefined vehicle-count or speed thresholds.}

\section{Apache Flink}
\subsection{Introduction and Architecture}
\placeholder{Explain Apache Flink and its stream-processing concepts,
execution model, data streams, operators, and windows, as applicable.}

\subsection{Real-Time Patient Data Processing}
\placeholder{Describe the synthetic patient vital-sign stream and the
processing flow used to identify readings outside the selected
demonstration thresholds.}

\subsection{Alert Generation and Time-Window Summary}
\placeholder{Explain how alerts are generated and how abnormal readings
are summarized over a specified time window. Clearly state that the
prototype is educational and is not a clinical diagnostic tool.}

\section{MongoDB}
\subsection{Introduction and Document Model}
\placeholder{Explain MongoDB, NoSQL databases, JSON-like/BSON documents,
collections, and the document-oriented data model.}

\subsection{Collections, Indexing, and CRUD Operations}
\placeholder{Describe the collection structure, indexes used, and
Create, Read, Update, and Delete operations implemented.}

\subsection{Proposed E-Commerce Product Catalog}
\placeholder{Explain how product details such as name, category, price,
brand, specifications, rating, and availability are represented and
queried.}

\section{Cloud-Based Big Data Analytics}
\subsection{Cloud Platform Overview}
\placeholder{Identify the selected platform (AWS, Microsoft Azure, or
Google Cloud) and introduce the storage and analytics services used.}

\subsection{Storage and Analytics Workflow}
\placeholder{Describe the selected managed storage and analytics services,
how the sample server-log dataset is uploaded or accessed, and how the
data is analyzed.}

\subsection{Log Analytics Design}
\placeholder{Explain the analysis of error frequencies, request volumes,
unusual spikes in failed requests, and time-based usage patterns. Include
the visualization approach and cost considerations.}

% =========================================================
% CHAPTER 4 - IMPLEMENTATION AND RESULTS
% =========================================================
\chapter{IMPLEMENTATION AND RESULTS}

\section{Implementation Environment}
\placeholder{Mention the operating system, hardware, software versions,
programming languages, tools, configuration, and execution environment
used for the activities.}

% ---------------- Activity 1 ----------------
\activityheading{1}{Apache Hive -- Supermarket Sales Analysis}

\subsection{Problem Statement}
A supermarket chain generates millions of sales transaction records
daily. Management wants to identify best-selling products, peak sales
hours, branch-wise revenue, and monthly sales trends to improve
inventory planning and business decisions.

\subsection{Self-Learning and Implementation Tasks}
\begin{enumerate}
    \item Study Hive architecture, HiveQL, and data warehousing concepts.
    \item Create Hive tables to store sales transaction data.
    \item Write queries to calculate total revenue, top-selling products,
    and branch-wise sales. Include peak-hour and monthly analysis where
    the dataset contains the required fields.
    \item Analyze the results and recommend inventory-management strategies.
\end{enumerate}

\subsection{Dataset and Table Design}
\placeholder{Describe the sales dataset, source, number of records, fields,
data types, and Hive database/table schema.}

\subsection{Implementation Procedure}
\begin{enumerate}
    \item Configure and start the required Hive environment.
    \item Create the database and sales table.
    \item Load or access the sales transaction dataset.
    \item Execute HiveQL queries for revenue, product, branch, and time-based analysis.
    \item Save query outputs for interpretation.
\end{enumerate}

\subsection{HiveQL Queries and Outputs}
\placeholder{Insert the HiveQL statements used. For each query, explain
its purpose and include the corresponding output.}

\begin{table}[H]
\centering
\caption{Hive Sales Analysis Results}
\begin{tabularx}{\textwidth}{|X|X|X|}
\hline
\textbf{Analysis} & \textbf{Query / Metric} & \textbf{Observation}\\
\hline
Total Revenue & [Query or metric] & [Observation]\\
\hline
Top-Selling Products & [Query or metric] & [Observation]\\
\hline
Branch-Wise Revenue & [Query or metric] & [Observation]\\
\hline
Peak Sales Hours & [Query or metric] & [Observation]\\
\hline
Monthly Sales Trends & [Query or metric] & [Observation]\\
\hline
\end{tabularx}
\end{table}

\subsection{Results and Business Insights}
\placeholder{Insert screenshots of table creation, data loading, queries,
and outputs. Explain the findings and inventory-management recommendations.}

% ---------------- Activity 2 ----------------
\clearpage
\activityheading{2}{Apache Kafka -- Real-Time Traffic Monitoring}

\subsection{Problem Statement}
A smart-city traffic control center receives continuous vehicle-count
and traffic-speed data from sensors installed at major junctions. It
needs to identify congestion-prone roads and generate timely alerts for
traffic operators.

\subsection{Self-Learning and Implementation Tasks}
\begin{enumerate}
    \item Study Kafka topics, producers, consumers, and brokers.
    \item Generate or simulate traffic sensor data and publish it to a Kafka topic.
    \item Develop a consumer application to process incoming traffic events.
    \item Identify congestion using predefined speed or vehicle-count thresholds
    and generate alerts.
\end{enumerate}

\subsection{Dataset and Kafka Configuration}
\placeholder{Describe the generated/simulated traffic data, broker
configuration, topic name, partitions, producer settings, and consumer
settings used.}

\subsection{Implementation Procedure}
\begin{enumerate}
    \item Set up and start the Kafka broker environment.
    \item Create the traffic-data topic.
    \item Develop or configure a producer to publish traffic events.
    \item Develop a consumer to receive and process events.
    \item Apply the selected congestion thresholds and generate alerts.
\end{enumerate}

\subsection{Producer and Consumer Code}
\placeholder{Insert the producer and consumer code, configuration
commands, and explanations.}

\subsection{Results and Congestion Alerts}
\begin{table}[H]
\centering
\caption{Sample Traffic Monitoring Results}
\begin{tabularx}{\textwidth}{|l|c|c|X|}
\hline
\textbf{Junction} & \textbf{Vehicle Count} & \textbf{Speed} & \textbf{Detected Status}\\
\hline
[Junction 1] & [Value] & [Value] & [Normal / Congested]\\
\hline
[Junction 2] & [Value] & [Value] & [Normal / Congested]\\
\hline
[Junction 3] & [Value] & [Value] & [Normal / Congested]\\
\hline
\end{tabularx}
\end{table}
\placeholder{Insert screenshots of topic creation, producer messages,
consumer output, and generated congestion alerts. Explain the observations.}

% ---------------- Activity 3 ----------------
\clearpage
\activityheading{3}{Apache Flink -- Real-Time Patient Vital-Sign Monitoring}

\subsection{Problem Statement}
A hospital monitors patients using connected medical devices that
continuously transmit heart rate, oxygen saturation, and body temperature
readings. The hospital wants to process readings in real time and flag
potentially abnormal values for clinical review.

\subsection{Self-Learning and Implementation Tasks}
\begin{enumerate}
    \item Study Flink architecture and stream-processing concepts.
    \item Simulate a continuous stream of patient vital-sign readings.
    \item Develop a Flink application to identify readings outside predefined
    demonstration thresholds.
    \item Generate alerts and summarize abnormal readings over a specified
    time window.
\end{enumerate}

\subsection{Dataset and Demonstration Thresholds}
\placeholder{Describe the synthetic data fields, data-generation method,
and demonstration thresholds used for heart rate, oxygen saturation,
and body temperature. State the source or rationale for the educational
thresholds.}

\subsection{Implementation Procedure}
\begin{enumerate}
    \item Configure the Flink execution environment.
    \item Generate or read the synthetic patient vital-sign stream.
    \item Define the stream-processing operations.
    \item Apply the selected demonstration thresholds.
    \item Generate alert records and time-window summaries.
\end{enumerate}

\subsection{Flink Program and Processing Logic}
\placeholder{Insert the Flink program, configuration, processing logic,
and explanation of the operators or windows used.}

\subsection{Results and Alert Summary}
\begin{table}[H]
\centering
\caption{Sample Patient Stream-Processing Results}
\begin{tabularx}{\textwidth}{|l|l|c|c|c|X|}
\hline
\textbf{Device/ID} & \textbf{Timestamp} & \textbf{Heart Rate} &
\textbf{SpO$_2$} & \textbf{Temperature} & \textbf{Status}\\
\hline
[ID 1] & [Time] & [Value] & [Value] & [Value] & [Normal / Flagged]\\
\hline
[ID 2] & [Time] & [Value] & [Value] & [Value] & [Normal / Flagged]\\
\hline
[ID 3] & [Time] & [Value] & [Value] & [Value] & [Normal / Flagged]\\
\hline
\end{tabularx}
\end{table}

\noindent\textbf{Important note:} This activity uses synthetic patient data
and educational thresholds. The prototype is intended only to demonstrate
stream processing and must not be treated as a clinical diagnostic tool.

\placeholder{Insert screenshots of the input stream, Flink execution,
flagged readings, alerts, and time-window summaries. Interpret the results.}

% ---------------- Activity 4 ----------------
\clearpage
\activityheading{4}{MongoDB -- E-Commerce Product Catalog}

\subsection{Problem Statement}
An e-commerce company manages millions of product records containing
different attributes, including product names, prices, brands,
specifications, ratings, and availability. It needs a flexible database
to manage and retrieve product information efficiently.

\subsection{Self-Learning and Implementation Tasks}
\begin{enumerate}
    \item Study MongoDB document storage, collections, indexing, and CRUD operations.
    \item Design a product catalog using JSON-like documents.
    \item Implement product insertion, category and price-range searches,
    stock updates, and deletion of discontinued products.
    \item Compare the flexibility of the document model with a relational database.
\end{enumerate}

\subsection{Document Structure and Collection Design}
\placeholder{Explain the product document structure, collection name,
fields, data types, and sample JSON-like documents.}

\subsection{Implementation Procedure}
\begin{enumerate}
    \item Create or connect to the MongoDB database.
    \item Create the product collection and insert sample product documents.
    \item Implement product retrieval and category/price-range searches.
    \item Update stock or availability details.
    \item Delete discontinued products.
    \item Create indexes where appropriate and observe query performance.
\end{enumerate}

\subsection{CRUD Queries and Outputs}
\placeholder{Insert MongoDB commands or application code for Create,
Read, Update, and Delete operations. Include category searches,
price-range queries, and stock updates.}

\subsection{Indexing and Query Performance}
\placeholder{Describe the indexes created, queries evaluated, and
observations about query performance. Do not report performance
improvements unless they were measured.}

\subsection{Document Database and Relational Database Comparison}
\begin{table}[H]
\centering
\caption{Comparison of MongoDB Document Model and Relational Database}
\begin{tabularx}{\textwidth}{|l|X|X|}
\hline
\textbf{Aspect} & \textbf{MongoDB} & \textbf{Relational Database}\\
\hline
Data Model & Document-oriented & Table-oriented\\
\hline
Schema & [Observation from implementation] & [Observation]\\
\hline
Representation of Product Attributes & [Observation] & [Observation]\\
\hline
Query Approach & [Observation] & [Observation]\\
\hline
Flexibility & [Observation] & [Observation]\\
\hline
\end{tabularx}
\end{table}

\subsection{Results and Observations}
\placeholder{Insert screenshots of database/collection creation,
document insertion, retrieval, updates, deletion, indexes, and query
outputs. Summarize the observations.}

% ---------------- Activity 5 ----------------
\clearpage
\activityheading{5}{Cloud-Based Big Data Analytics -- Application Log Analysis}

\subsection{Problem Statement}
An online service provider generates gigabytes of application and
web-server logs daily. The operations team wants to identify frequent
errors, analyze request patterns, detect unusual spikes in failed
requests, and estimate peak usage periods without maintaining
large-scale infrastructure.

\subsection{Self-Learning and Implementation Tasks}
\begin{enumerate}
    \item Explore a cloud-based Big Data platform.
    \item Select suitable managed storage and analytics services.
    \item Upload a sample server-log dataset and analyze error frequencies,
    request volumes, and time-based usage patterns.
    \item Document the architecture, processing workflow, findings, and
    estimated resource costs.
\end{enumerate}

\subsection{Selected Cloud Platform and Services}
\placeholder{Select and identify the platform used: AWS, Microsoft Azure,
or Google Cloud. List the exact storage and analytics services used.
If cloud access is unavailable, describe the local simulation or free-tier
approach adopted.}

\begin{table}[H]
\centering
\caption{Selected Cloud Services}
\begin{tabularx}{\textwidth}{|l|X|}
\hline
\textbf{Item} & \textbf{Details}\\
\hline
Cloud Platform & [AWS / Microsoft Azure / Google Cloud]\\
\hline
Storage Service & [Selected service]\\
\hline
Analytics / Query Service & [Selected service]\\
\hline
Visualization Service / Tool & [Selected service or tool]\\
\hline
Dataset & [Sample log dataset and source]\\
\hline
\end{tabularx}
\end{table}

\subsection{Dataset and Log Format}
\placeholder{Describe the sample server-log dataset, its format, fields,
size, source, and any preprocessing performed.}

\subsection{Architecture and Processing Workflow}
\placeholder{Explain the data flow from log upload/collection to storage,
querying, analysis, and visualization. Insert the architecture diagram
if available.}

\begin{figure}[H]
    \centering
    \fbox{\parbox[c][5cm][c]{0.85\textwidth}{\centering
    Insert Cloud-Based Log Analytics Workflow Diagram}}
    \caption{Cloud-Based Big Data Log Analytics Workflow}
    \label{fig:cloud_log_workflow}
\end{figure}

\subsection{Implementation Procedure}
\begin{enumerate}
    \item Select and configure the cloud platform or local simulation.
    \item Create or configure the required storage resource.
    \item Upload or access the sample application/web-server logs.
    \item Configure the analytics service and execute the required queries.
    \item Analyze error frequencies, request volumes, failed-request spikes,
    and time-based usage patterns.
    \item Generate visualizations and document estimated resource costs.
\end{enumerate}

\subsection{Queries, Analysis, and Visualizations}
\placeholder{Insert the queries, scripts, or processing steps used.
Include output tables or visualizations for frequent errors, request
volumes, failed requests, and peak usage periods.}

\subsection{Cost Considerations}
\placeholder{Document the services and resource usage considered,
estimated costs if available, and whether a free tier or local simulation
was used. Do not include credentials, access keys, or confidential
account information.}

\subsection{Results and Observations}
\placeholder{Insert screenshots of storage, dataset upload, analytics
queries, outputs, and visualizations. Interpret the findings and explain
how they could help the operations team.}

% =========================================================
% CHAPTER 5 - OVERALL ANALYSIS
% =========================================================
\clearpage
\chapter{OVERALL ANALYSIS AND DISCUSSION}

\section{Summary of Activity Outcomes}
\begin{table}[H]
\centering
\caption{Summary of Technology Implementation Outcomes}
\begin{tabularx}{\textwidth}{|l|X|X|}
\hline
\textbf{Technology} & \textbf{Implementation Output} & \textbf{Key Observation}\\
\hline
Apache Hive & Sales analysis queries and report & [Observation]\\
\hline
Apache Kafka & Producer-consumer prototype and traffic alerts & [Observation]\\
\hline
Apache Flink & Stream-processing prototype and alert summaries & [Observation]\\
\hline
MongoDB & Product catalog with CRUD operations and queries & [Observation]\\
\hline
Cloud Analytics & Log analytics prototype/report and visualizations & [Observation]\\
\hline
\end{tabularx}
\end{table}

\section{Comparative Discussion}
\placeholder{Compare the roles of the five technologies. Discuss the
differences between data warehousing, event streaming, stream processing,
document databases, and cloud-based analytics based on the activities
performed.}

\section{Limitations}
\begin{itemize}
    \item [Mention dataset size or availability limitations.]
    \item [Mention configuration, hardware, or cloud-access limitations.]
    \item [Mention simulation assumptions and threshold limitations.]
    \item [Mention any implementation features not completed.]
\end{itemize}

% =========================================================
% CHAPTER 6 - CONCLUSION AND FUTURE WORK
% =========================================================
\chapter{CONCLUSION AND FUTURE WORK}

\section{Conclusion}
\placeholder{Summarize the five activities, technologies studied,
implementations completed, results obtained, and the overall learning
outcomes. Ensure that the conclusion reflects the actual work performed.}

\section{Future Work}
\begin{itemize}
    \item Extend the Hive analysis using larger and more varied sales datasets.
    \item Enhance Kafka traffic monitoring with additional sensor sources and
    more detailed event-processing logic.
    \item Explore more advanced Flink windowing and event-time processing.
    \item Extend the MongoDB product catalog with additional search and
    indexing strategies.
    \item Explore additional cloud analytics capabilities and dashboards,
    while monitoring resource usage and cost.
\end{itemize}

% =========================================================
% REFERENCES
% =========================================================
\clearpage
\addcontentsline{toc}{chapter}{References}
\begin{thebibliography}{99}

\bibitem{hive}
Apache Software Foundation, ``Apache Hive Documentation,''
\url{https://hive.apache.org/}.

\bibitem{kafka}
Apache Software Foundation, ``Apache Kafka Documentation,''
\url{https://kafka.apache.org/documentation/}.

\bibitem{flink}
Apache Software Foundation, ``Apache Flink Documentation,''
\url{https://flink.apache.org/}.

\bibitem{mongodb}
MongoDB, ``MongoDB Documentation,''
\url{https://www.mongodb.com/docs/}.

\bibitem{aws}
Amazon Web Services, ``AWS Documentation,''
\url{https://docs.aws.amazon.com/}.

\bibitem{azure}
Microsoft, ``Microsoft Azure Documentation,''
\url{https://learn.microsoft.com/azure/}.

\bibitem{gcp}
Google Cloud, ``Google Cloud Documentation,''
\url{https://cloud.google.com/docs}.

\bibitem{other}
[Add textbooks, research papers, articles, datasets, and other resources
actually referred to during the work.]

\end{thebibliography}

% =========================================================
% APPENDIX A - ASSESSMENT RUBRIC
% =========================================================
\appendix
\chapter{COMMON ASSESSMENT RUBRIC}

The activities are evaluated using the common assessment rubric below.
Each of the five activities carries 4 marks, for a total of 20 marks.

\begin{table}[H]
\centering
\caption{Common Assessment Rubric}
\begin{tabularx}{\textwidth}{|c|X|c|c|}
\hline
\textbf{Sl. No.} & \textbf{Evaluation Criterion} &
\textbf{Marks per Activity} & \textbf{Total for 5 Activities}\\
\hline
1 & Understanding of the problem and technology & 1 & 5\\
\hline
2 & Solution design and implementation & 1.5 & 7.5\\
\hline
3 & Results, analysis, and interpretation & 1 & 5\\
\hline
4 & Documentation and conclusions & 0.5 & 2.5\\
\hline
\multicolumn{2}{|r|}{\textbf{TOTAL}} & \textbf{4} & \textbf{20}\\
\hline
\end{tabularx}
\end{table}

% =========================================================
% APPENDIX B - SUBMISSION REQUIREMENTS
% =========================================================
\chapter{REPORT SUBMISSION REQUIREMENTS}

The final consolidated report should include the following components
for the activities, as applicable.

\begin{longtable}{|c|p{3.3cm}|p{8.0cm}|}
\caption{Mandatory Report Components}\\
\hline
\textbf{Sl. No.} & \textbf{Component} & \textbf{Description}\\
\hline
\endfirsthead
\hline
\textbf{Sl. No.} & \textbf{Component} & \textbf{Description}\\
\hline
\endhead
1 & Problem Definition & Clearly explain the real-world problem and objectives.\\
\hline
2 & Technology Study & Describe the selected technology, architecture, and relevant concepts.\\
\hline
3 & Implementation & Include the procedure, queries or code, and configuration details as applicable.\\
\hline
4 & Results & Provide screenshots, sample outputs, tables, or visualizations.\\
\hline
5 & Analysis & Interpret the results and explain how the solution addresses the problem.\\
\hline
6 & Conclusion & Summarize findings, limitations, and possible improvements.\\
\hline
\end{longtable}

\section*{Final Submission Checklist}
\addcontentsline{toc}{section}{Final Submission Checklist}
\begin{itemize}
    \item[$\square$] Problem definition and objectives included for all five activities.
    \item[$\square$] Technology architecture and concepts explained.
    \item[$\square$] Implementation procedures, queries, code, and configurations included.
    \item[$\square$] Screenshots, sample outputs, tables, and visualizations included.
    \item[$\square$] Results analyzed and interpreted.
    \item[$\square$] Limitations and conclusions documented.
    \item[$\square$] References and dataset sources listed.
    \item[$\square$] Cloud costs or local/free-tier simulation details documented, where applicable.
    \item[$\square$] Synthetic patient data and educational-only limitation clearly stated for the Flink activity.
\end{itemize}

\end{document}
